\newpage
\section{Acta Número 08}
\label{sec:acta08}

\noindent \textbf{Datos de la Sesión}
\begin{itemize}[noitemsep]
    \item \textbf{Modalidad:} Virtual - Google Meet.
    \item \textbf{Hora de Inicio:} 18:00
    \item \textbf{Hora de Fin:} 20:00
\end{itemize}

\noindent \textbf{Orden del Día}
\begin{itemize}[noitemsep]
    \item \textbf{Asistencia.}
    \item \textbf{Desglose de Épicas (Epic 1 y Epic 2):} Descomposición de las épicas principales en Historias de Usuario para el Sprint 1.
    \item \textbf{Establecimiento de la \textit{Definition of Ready} (DoR):} Aprobación de los criterios y condiciones previas necesarias para que una Historia de Usuario ingrese al Sprint.
    \item \textbf{Estimación de Esfuerzo mediante \textit{Planning Poker}:} Evaluación de la complejidad de las Historias de Usuario utilizando Puntos de Historia (Serie de Fibonacci).
    \item \textbf{Consolidación de Backlogs:} Creación del \textit{Product Backlog}, selección de PBIs y conformación del \textit{Sprint Backlog} para el Sprint 1.
    \item \textbf{Firma y aprobación de los presentes.}
\end{itemize}

\noindent \textbf{Desarrolladores (Socios Fundadores)}
\begin{enumerate}[noitemsep]
    \item Pardo Pozo Juan Diego (Jefe de Proyecto).
    \item Arnez Jaimes Mauricio.
    \item Quispe Flores Camila Belen.
    \item Ariñez Bautista Jim Gabriel.
    \item Medina Rodríguez Mateo Jhoustin.
    \item Molina Maldonado Tomas Manuel.
    \item Zeballos Crespo Jonathan.
    \item Gutierrez Guzmán Angheline.
\end{enumerate}

\noindent \textbf{Fecha de la sesión:}\par
Martes 31 de marzo de 2026

\newpage
\subsection{Desarrollo del Acta}

\noindent \textbf{Asistencia}\par
Se procedió a la toma de asistencia a la hora acordada, corroborando la presencia en sesión de los 8 socios fundadores de la grupo-empresa Byte Busters S.R.L. Al contar con la totalidad de los miembros, se estableció el quórum necesario para iniciar la toma de decisiones y el trabajo de planificación.

\vspace*{0.5cm}
\noindent\textbf{Desglose de Épicas (Epic 1 y Epic 2)}\par
El equipo técnico llevó a cabo el análisis y descomposición de la Epic 1 y la Epic 2. Se dividieron estos requerimientos de alto nivel en Historias de Usuario granulares y manejables, garantizando que cada una posea un alcance bien definido y aporte valor tangible al incremento del producto proyectado para el Sprint 1.

\vspace*{0.5cm}
\noindent \textbf{Establecimiento de la \textit{Definition of Ready} (DoR)}\par
Se debatió y consensuó la \textit{Definition of Ready} (DoR) estándar para el equipo. Se estableció que ninguna Historia de Usuario podrá ingresar a un Sprint si no cumple con los siguientes parámetros:
\begin{itemize}[noitemsep]
    \item \textbf{Claridad:} Comprensión absoluta y compartida sobre qué se debe desarrollar y su propósito comercial o técnico.
    \item \textbf{Criterios de Aceptación:} Definición exacta de las condiciones y pruebas que la funcionalidad debe superar para considerarse finalizada.
    \item \textbf{Estimación:} Complejidad y esfuerzo evaluados y asignados por el equipo de desarrollo.
    \item \textbf{Independencia:} Capacidad de ser desarrollada y testeada sin bloqueos generados por otras dependencias externas u otros equipos.
    \item \textbf{Tamaño Adecuado:} Alcance lo suficientemente acotado como para ser completado íntegramente dentro de la iteración de un solo Sprint.
    \item \textbf{Recursos Listos:} Disponibilidad total de accesos, credenciales, diseños de interfaz (UI) y esquemas de datos necesarios para iniciar la codificación de inmediato.
\end{itemize}

\vspace*{0.5cm}
\noindent \textbf{Estimación de Esfuerzo mediante \textit{Planning Poker}}\par
Con las Historias de Usuario definidas y la DoR aprobada, el equipo procedió a estimar el esfuerzo de las tareas correspondientes a la Epic 1 y 2. Se utilizó la técnica de \textit{Planning Poker} basándose en la sucesión de Fibonacci. Se enfatizó y respetó el principio de estimar basándose en el nivel de complejidad, riesgo e incertidumbre técnico (\textit{Story Points}), descartando explícitamente la estimación tradicional basada en horas de trabajo.

\vspace*{0.5cm}
\noindent \textbf{Consolidación del \textit{Product Backlog} y \textit{Sprint Backlog}}\par
Se estructuró oficialmente el \textit{Product Backlog} y se identificaron los Elementos del Backlog del Producto (PBIs) de mayor prioridad. A partir de estos, se construyó y documentó el \textit{Sprint Backlog} que regirá el ciclo de trabajo del Sprint 1. Se estableció el día de hoy, 31 de marzo, a la medianoche (00:00 hrs) como fecha y hora límite improrrogable para la entrega final y registro de estos artefactos ágiles en las plataformas de gestión del equipo.

\vspace*{0.5cm}
\noindent \textbf{Aprobación Oficial}\par
Habiendo estructurado el plan de iteración, definido los requerimientos de calidad para el inicio del trabajo (DoR) y consolidado los backlogs, los 8 socios fundadores aprobaron por unanimidad las planificaciones realizadas durante la sesión, oficializando el arranque técnico del Sprint 1.

\vspace*{0.5cm}
\noindent \textbf{Firmas y aprobación de los integrantes}
\begin{center}
    \noindent
    \setlength{\tabcolsep}{0pt}
    
    \begin{tabular}{ccc}
        \espacioFirma{assets/img/firmas/mauricio.jpg}{Mauricio Jaimes Arnez}{13751221} & 
        \espacioFirma{assets/img/firmas/camila.jpg}{Camila Belen Quispe Flores}{13097244} &
        \espacioFirma{assets/img/firmas/jim.jpg}{Jim Gabriel Ariñez Bautista}{9517642} \\[1.0cm]
        
        \espacioFirma{assets/img/firmas/mateo.jpg}{Mateo Medina Rodríguez}{9453809} &
        \espacioFirma{assets/img/firmas/tomas.jpg}{Tomas Molina Maldonado}{8051701} &
        \espacioFirma{assets/img/firmas/jonathan.jpg}{Jonathan Zeballos Crespo}{13067494} \\[1.0cm]
    \end{tabular}

    \vspace{0.2cm} 
    \begin{tabular}{cc}
        \espacioFirma{assets/img/firmas/diego.jpg}{Juan Diego Pardo Pozo}{12747783} & 
        \espacioFirma{assets/img/firmas/angheline.jpg}{Angheline Gutierrez Guzmán}{13751815}
    \end{tabular}
\end{center}

\vspace{0.5cm}
\noindent \textbf{Fecha de última revisión:} 31 de marzo de 2026